\documentclass[a4paper,14pt]{extarticle}

\usepackage{polyglossia}
\setdefaultlanguage{russian}

\usepackage{amsmath, amssymb}
\usepackage{graphicx}
\usepackage{geometry}
\usepackage{setspace}
\usepackage{indentfirst}
\usepackage{fontspec}
\usepackage{tabularx}
\usepackage{colortbl}
\usepackage{pgfplots}

\pgfplotsset{compat=1.18}

\definecolor{lightgray}{gray}{0.9}

\setmainfont{Times New Roman}

\geometry{left=2cm,right=2cm,top=2cm,bottom=2cm}
\onehalfspacing
\setlength{\parindent}{1.25cm}

\begin{document}

\begin{titlepage}
\begin{center}

\textbf{Министерство образования Республики Беларусь}

\textbf{Белорусский государственный университет}

Факультет радиофизики и компьютерных технологий

\vspace{5cm}

\textbf{РЕФЕРАТ}

\vspace{0.5cm}

на тему

\textbf{«ПАРАМАГНЕТИЗМ»}

\end{center}

\vfill

\begin{flushright}

Выполнил студент гр. 7 \\

Нацевский Матвей

\vspace{0.5cm}

Преподаватель Садов П. В.

\end{flushright}

\vfill

\begin{center}
Минск, 2026
\end{center}

\end{titlepage}

\tableofcontents
\newpage

\section{Введение}

Изучение магнитных свойств материи традиционно считается одним из наиболее фундаментальных разделов физики конденсированного состояния. Понимание механизмов взаимодействия вещества с внешним магнитным полем позволяет не только раскрыть внутреннюю структуру атомов и молекул, но и создавать принципиально новые устройства в области электроники, медицины и квантовых вычислений. Парамагнетизм, как особое состояние вещества, характеризующееся способностью намагничиваться вдоль внешнего поля, представляет собой уникальный объект исследования, где макроскопические термодинамические параметры напрямую определяются квантовыми характеристиками микрочастиц.

Актуальность данной темы в контексте современного научного знания обусловлена стремительным развитием нанотехнологий и молекулярного магнетизма. Если на заре изучения магнетизма в XIX веке ученые оперировали макроскопическими телами, то сегодня фокус исследователей сместился к отдельным молекулам и наночастицам, проявляющим суперпарамагнитные свойства.

Целью данного реферата является систематическое изложение теоретических основ парамагнетизма, анализ его микроскопической природы и обзор современных прикладных аспектов.

Для достижения цели необходимо решить следующие задачи:

\begin{itemize}
\item раскрыть физическую природу магнитных моментов микрочастиц;
\item проанализировать классические теории Кюри и Ланжевена;
\item описать квантово-механическую модель парамагнетизма;
\item рассмотреть специфические формы парамагнетизма;
\item изучить методы исследования парамагнитных веществ;
\item оценить практическое применение парамагнетизма.
\end{itemize}

\section{Природа парамагнетизма и магнитные моменты микрочастиц}

Физическая сущность парамагнетизма заключается в наличии у структурных единиц вещества — атомов, ионов или молекул — постоянных магнитных моментов. В отличие от диамагнетиков, где магнитный момент индуцируется внешним полем и направлен против него, в парамагнетиках поле лишь ориентирует уже существующие микромагниты.

Магнитный момент атома складывается из нескольких квантовых компонент. Основными носителями магнетизма являются электроны. Их движение можно разделить на орбитальное и спиновое. Орбитальное движение электрона вокруг ядра создаёт магнитный момент, аналогичный току в замкнутом контуре. Спиновый момент связан с внутренним моментом импульса электрона.

Для характеристики магнитных моментов используется магнетон Бора

\[
\mu_B = \frac{e\hbar}{2m_e}.
\]

Парамагнетизм проявляется в системах, где присутствуют неспаренные электроны. К таким веществам относятся:

\begin{itemize}
\item ионы переходных металлов;
\item свободные радикалы;
\item молекулярный кислород $O_2$;
\item некоторые металлы с электронами проводимости.
\end{itemize}

Энергия взаимодействия магнитного момента с внешним полем определяется выражением

\[
U = -\vec{\mu}\cdot\vec{B}.
\]

\subsection*{Магнитная восприимчивость некоторых парамагнетиков}

\begin{center}
\begin{tabular}{|c|c|c|}
\hline
Материал & Символ & $\chi_m \times 10^{-5}$ \\
\hline
Вольфрам & W & 7.8 \\
Цезий & Cs & 6.1 \\
Натрий & Na & 3.72 \\
Алюминий & Al & 3.2 \\
Литий & Li & 2.4 \\
Магний & Mg & 2.2 \\
\hline
\end{tabular}
\end{center}

\section{Классическая теория парамагнетизма}

\subsection{Закон Кюри}

Пьер Кюри установил, что магнитная восприимчивость парамагнетиков обратно пропорциональна температуре:

\[
\chi = \frac{C}{T}.
\]

Здесь $C$ — постоянная Кюри.

Она выражается как

\[
C = \frac{\mu_0 n \mu^2}{3k_B}.
\]

\subsection{Теория Ланжевена}

Ланжевен рассмотрел парамагнетик как систему частиц с магнитным моментом $\mu$.

Средняя намагниченность определяется функцией Ланжевена:

\[
L(x)=\coth x - \frac{1}{x},
\]

где

\[
x=\frac{\mu B}{k_B T}.
\]

В слабых полях

\[
L(x)\approx \frac{x}{3}.
\]

Это приводит к закону Кюри.

\section{Квантово-механическая интерпретация}

Квантовая теория вводит пространственное квантование магнитного момента.

Проекция момента определяется выражением

\[
\mu_z = - g \mu_B m_J.
\]

Намагниченность описывается функцией Бриллюэна

\[
B_J(x)=
\frac{2J+1}{2J}\coth\left(\frac{2J+1}{2J}x\right)
-
\frac{1}{2J}\coth\left(\frac{x}{2J}\right).
\]

\section{Особенности парамагнетизма в конденсированных средах}

\subsection{Парамагнетизм Паули}

В металлах магнитная восприимчивость почти не зависит от температуры и определяется плотностью состояний на уровне Ферми.

\subsection{Парамагнетизм Ван Флека}

В некоторых ионах основной уровень имеет $J=0$. Магнитный момент возникает за счёт смешивания возбужденных состояний внешним полем.

\subsection{Суперпарамагнетизм}

В наночастицах размером менее 20–30 нм каждая частица ведёт себя как один магнитный домен. Их поведение напоминает парамагнетик, но с огромной восприимчивостью.

\section{Экспериментальные методы исследования}

Наиболее точным методом является электронный парамагнитный резонанс (ЭПР), открытый Е.~К.~Завойским.

Резонанс возникает при условии

\[
h\nu = g\mu_B B.
\]

\section{Прикладное значение}

Парамагнетизм используется:

\begin{itemize}
\item для получения сверхнизких температур методом адиабатического размагничивания;
\item в медицинской диагностике (контрастные вещества для МРТ);
\item в исследованиях структуры веществ.
\end{itemize}

\section{Современные направления}

Перспективной областью является молекулярный магнетизм. Одномолекулярные магниты способны сохранять намагниченность на уровне отдельной молекулы и рассматриваются как возможная основа для квантовых вычислений.

\section{Заключение}

Парамагнетизм представляет собой важное физическое явление, связывающее квантовую природу микрочастиц с макроскопическими магнитными свойствами вещества. Современные исследования открывают новые области применения парамагнитных материалов — от нанотехнологий до квантовых вычислений.

\section*{Список литературы}

\begin{enumerate}
\item Альтшулер С. А., Козырев Б. М. Электронный парамагнитный резонанс соединений элементов промежуточных групп. — М.: Наука, 1972.
\item Ашкрофт Н., Мермин Н. Физика твердого тела. — М.: Мир, 1979.
\item Белов К. П. Магнитные превращения. — М.: Физматгиз, 1959.
\item Блейкмор Дж. Физика твердого тела. — М.: Мир, 1988.
\item Вонсовский С. В. Магнетизм. — М.: Наука, 1971.
\item Кириченко Н. А. Электричество и магнетизм. — М.: МФТИ, 2011.
\item Киттель Ч. Введение в физику твердого тела. — М.: Наука, 1978.
\item Сивухин Д. В. Общий курс физики. Т. III. — М.: Физматлит, 2004.
\item Тамм И. Е. Основы теории электричества. — М.: Физматлит, 2003.
\item Фейнман Р. Фейнмановские лекции по физике. — М.: Мир, 1977.
\end{enumerate}

\end{document}